\begin{annexitem}{annex:p15.rat53.rank-laws}{thm:p15-rat53-rank-laws}{Event laws, world representation and threshold beliefs}
\begin{proof}
If \(A\subseteq B\), then \(A\lor B=B\), so the minimum law gives
\(\kappa(B)\leq\kappa(A)\).  Applying it to \(A\lor\neg A=W\)
gives the negation law.  For two believed events at threshold \(z\),
\[
 \kappa(\neg(A\land B))=
 \min\{\kappa(\neg A),\kappa(\neg B)\}>z.
\]
Induction gives finite conjunction closure.  The contradiction is not
believed, since its negation has rank zero.
De Morgan's law gives the positive version's conjunction law and endpoint
values.  Since at least one of \(\kappa(A)\), \(\kappa(\neg A)\)
is zero, its signed difference determines the nonzero one by its sign and
absolute value; this proves both inverse formulas, including infinite values.

On finite \(W\), repeated binary minimum preservation on the union of
singletons gives the world-minimum formula.  Normalization supplies at
least one world of rank zero, so \(M_z\ne\varnothing\).
An event's complement has rank greater than \(z\) exactly when every
world of rank at most \(z\) is in the event.  This proves
\(\mathcal B_z=\{A:M_z\subseteq A\}\).  Its intersection is \(M_z\)
because \(M_z\) itself is one of these events.  This proves both the
nonempty intersection and closure under the source's full intersection-based
consequence.  The same argument works on any space with an attained
world-minimum representation; the set of worlds of rank at most \(z\)
still contains a zero-rank world.

For the converse write \(\beta(A)=\kappa(\neg A)\).
Upward closure of all threshold sets implies \(A\subseteq B\Rightarrow
\beta(A)\leq\beta(B)\): if the first value were greater, the finite
threshold \(z=\beta(B)\) would witness failure of upward closure.  If
\(\beta(B)=\infty\) that strict inequality is impossible.
Hence \(\beta(A\land B)\leq\min\{\beta(A),\beta(B)\}\).
If the inequality were strict, its left side would be finite; choose it
as the threshold.  Both conjuncts would be accepted and their conjunction
would not, contradicting intersection closure.  Thus equality holds;
taking complements gives the binary minimum axiom.  Normalization was a
separate hypothesis, not manufactured by this argument.

For the infinite counterexample, a union is finite exactly when both of its
two members are finite.  This proves the binary minimum axiom of the
displayed event function.  Its believed events are exactly the cofinite sets.
Every point is omitted by one such set, so the full intersection is empty,
although every finite intersection is cofinite and nonempty.  No world-minimum
representation is possible: all singleton ranks are infinite but \(W\)
has rank zero.
\end{proof}
\end{annexitem}

\begin{annexitem}{annex:p15.rat53.conditional-laws}{thm:p15-rat53-conditional-laws}{Conditional ranking and the exact converse hypotheses}
\begin{proof}
Subset antitonicity gives \(\kappa(A\land B)\geq\kappa(A)\), so the
conditional is nonnegative.  Its whole-space value is zero and its
empty-space value infinite.  Subtracting the same finite \(\kappa(A)\)
from the two branches of the minimum law preserves minimum, giving the
conditional disjunction law.  In particular the partition of \(A\) into
\(A\land B\) and \(A\land\neg B\) has minimum rank \(\kappa(A)\);
at least one of the two conditional ranks is zero.  The conjunction identity
is the defining subtraction rewritten, with only finite subtraction.

Suppose \(A\subseteq B\) and the two stated ranks are finite.
Then \(\kappa(\neg A)\leq\kappa(\neg B)<\infty\).
Conditional on \(A\), \(B\) has negative rank zero and its negation
infinite, so \(\tau(B\mid A)=\infty\).
Conditional on \(\neg A\), the negative rank of \(\neg B\) is finite,
because \(\neg B\subseteq\neg A\); the two-sided rank of \(B\) is
therefore finite or negative infinity, but not positive infinity.
The strict reason comparison follows.  Dropping the finite-rank conditions
would invoke an undefined subtraction or compare certainty in both branches.

To justify the qualified alternative axiomatization, take normalization,
subset antitonicity, and the conditional negation law for each finite-rank
condition.  Put \(C=A\lor B\).  If \(\kappa(C)=\infty\), its subsets
both have infinite rank, so the minimum law holds.  Otherwise the
conditional law for the partition \(A\), \(C\setminus A\) says
\(\kappa(C)=\min\{\kappa(A),\kappa(C\setminus A)\}\).
Since \(C\setminus A\subseteq B\), antitonicity gives a lower bound
\(\min\{\kappa(A),\kappa(B)\}\) for this expression, while
\(A,B\subseteq C\) give the reverse inequality.  Thus the binary
minimum law follows.

Without antitonicity this reverse claim fails even on three worlds: assign
zero to each singleton and to \(W\), and infinity to every two-element
set and to the empty set.  Each finite-rank condition is either a singleton
or \(W\); partitioning it always gives at least one zero-rank part.
The conditional negation law therefore holds wherever its subtraction is
defined.  But the union of two zero-ranked singletons has infinite rank.
This is why conditional consistency alone is not used as an unqualified
replacement for the original axioms.
\end{proof}
\end{annexitem}

\begin{annexitem}{annex:p15.rat53.update}{thm:p15-rat53-update}{The uniquely specified two-cell dynamics and both Tweety steps}
\begin{proof}
Both conditional functions are rankings by the previous result.
Adding constants \(a,b\) and taking their pointwise minimum preserves
binary minimum preservation; explicitly the four terms for a union can be
regrouped as the minimum of the two posterior event ranks.
The empty event has rank infinity, and the whole space has rank
\(\min(a,b)=0\).  For \(E\), the first conditional rank is zero and
the second infinite, so its new rank is \(a\); the complement similarly
has new rank \(b\).
For any \(A\), the rank of \(A\land E\) under the update is
\(\kappa(A\mid E)+a\), since its complement-cell branch is impossible.
Subtracting its finite new cell rank \(a\) recovers the old conditional
\(\kappa(A\mid E)\).  The complement case is identical.

Conversely any posterior with those cell ranks and conditional functions
must give \(A\land E\) and \(A\land\neg E\) those two values.
Their union is \(A\), so its binary minimum law forces precisely the
displayed update formula.  This establishes uniqueness rather than merely
defining a named operation.  Applying it to a singleton world gives the
two piecewise offsets.  If a prior cell has infinite rank, its subtraction
is not defined; the proof does not cancel \(\infty-\infty\).

Hard conditioning uses just the previous conditional ranking function.
Its complement-evidence event intersects the condition in the empty set,
so its new rank is infinite.  This is different from the finite-firmness
case, and in particular can obstruct later conditioning on that complement.

For the eight source worlds in table order, the prior bird and nonbird
minimum ranks are both zero.  The first update therefore adds two exactly
to the two nonbird columns, producing the source's second table.
Its penguin minimum is one, at the nonflying penguin bird, and its
nonpenguin minimum remains zero.  The second update subtracts one from
both penguin columns and adds one to both nonpenguin columns, giving
exactly the third table.  The zero sets stated in the body follow by
inspection of the resulting eight values.

The conditional flight calculations use the prior minima:
\(\kappa(B)=0\), \(\kappa(B\land\neg F)=1\),
\(\kappa(B\land\neg P)=0\),
\(\kappa(B\land\neg P\land\neg F)=2\),
\(\kappa(P)=1\), \(\kappa(P\land F)=4\).
The three differences are respectively one, two and three.
The material penguin-to-flight implication is believed because its
negation \(P\land\neg F\) has rank one; the genuine conditional
instead disfavors flight by three.  This proves the source example's
claimed contrast without equating the two conditionals.
\end{proof}
\end{annexitem}

\begin{annexitem}{annex:p15.rat53.possibility}{thm:p15-rat53-possibility}{Static transforms, range restrictions and the probability-exponent bridge}
\begin{proof}
Exponentiation by \(2^{-x}\) reverses order on finite ranks, with infinity
sent to zero.  Thus the negative-rank minimum law becomes the possibility
maximum law, and normalization becomes \(\Pi(W)=1\),
\(\Pi(\varnothing)=0\).  Complementation followed by one minus gives
the claimed necessity minimum law and its endpoint values.
The inverse \(-\log_2\Pi(A)\) is an integer or infinity exactly on the
specified grid; for example a possibility value \(0.7\) cannot be
identified with an integer rank by this fixed transform.
For finite \(\kappa(A)\), exponentiating the defining subtraction gives
\(2^{-\kappa(B\mid A)}=\Pi(A\land B)/\Pi(A)\); its denominator
is positive.  The theorem does not prescribe an ordinal conditioning rule
with a different normalization convention.

For the finite regular probability family, write \(m=\min_{w\in A}r(w)\).
The numerator of \(P_\varepsilon(A)\) is
\(\varepsilon^m(c_A+o(1))\), where \(c_A\) is the positive number
of rank-\(m\) worlds in \(A\).  The denominator is
\(c_W+o(1)\), with \(c_W>0\) the number of zero-rank worlds.
Hence
\(\log P_\varepsilon(A)=m\log\varepsilon+O(1)\), and division by
\(\log\varepsilon\to-\infty\) gives the claimed limiting exponent.
Normalization, finite \(W\), finite world ranks and a nonempty event are
used explicitly.  This asymptotic family is not a unique ordinary
probability attached to \(\kappa\).

Finally, for the ten-outcome threshold counterexample each excluded-outcome
event has mass nine tenths and their intersection is empty.  This verifies
the concrete failure of probabilistic threshold closure; it does not assume
that no more restricted probabilistic account of outright belief is possible.
\end{proof}
\end{annexitem}

\begin{annexitem}{annex:p15-w02b-p15-039-reason-laws}{thm:p15-w02b-p15-039-reason-laws}{The cross-sum form of the reason relation}
\begin{proof}
The four marginal ranks are \(\kappa(A)=\min\{a,b\}\),
\(\kappa(\neg A)=\min\{c,d\}\), \(\kappa(B)=\min\{a,c\}\) and
\(\kappa(\neg B)=\min\{b,d\}\), by the binary minimum law, and all four
are finite by hypothesis.  By the definition of conditional ranks,
\(\kappa(B\mid A)=a-\kappa(A)\) and \(\kappa(\neg B\mid A)=b-\kappa(A)\),
with \(\kappa(A)\) finite.  Since \(\min\{a,b\}\) is finite, at most one of
\(a,b\) is infinite.  Hence \(\tau(B\mid A)=b-a\) when both are finite,
\(+\infty\) when \(b=\infty\), and \(-\infty\) when \(a=\infty\); no
\(\infty-\infty\) occurs.  In the same way \(\tau(B\mid\neg A)\) is
\(d-c\), \(+\infty\) or \(-\infty\).

\emph{Item 1.}  The cell \(a\) shares the row \(A\) with \(b\) and the
column \(B\) with \(c\).  If \(a=\infty\), finiteness of \(\min\{a,b\}\)
and \(\min\{a,c\}\) forces \(b,c<\infty\).  If \(d=\infty\), finiteness
of \(\min\{c,d\}\) and \(\min\{b,d\}\) forces \(b,c<\infty\).  So
\(a+d=\infty\) implies \(b+c<\infty\), and symmetrically \(b+c=\infty\)
implies \(a+d<\infty\).  There are three cases.
If all four cells are finite, \(\tau(B\mid A)>\tau(B\mid\neg A)\) reads
\(b-a>d-c\), which is \(a+d<b+c\).
If \(a+d=\infty\), then \(b,c\) are finite.  Either \(a=\infty\), and
\(\tau(B\mid A)=-\infty\) exceeds nothing; or \(a<\infty=d\), and
\(\tau(B\mid A)=b-a\) is finite while \(\tau(B\mid\neg A)=+\infty\).  In
both subcases \(A\) is not a reason for \(B\), and \(a+d<b+c\) is false.
If \(b+c=\infty\), then \(a,d\) are finite.  Either \(b=\infty\), and
\(\tau(B\mid A)=+\infty\) while \(\tau(B\mid\neg A)\) is finite or
\(-\infty\); or \(b<\infty=c\), and \(\tau(B\mid A)\) is finite while
\(\tau(B\mid\neg A)=-\infty\).  In both subcases \(A\) is a reason for
\(B\), and \(a+d<b+c\) is true.
Exchanging the roles of \(A\) and \(B\) exchanges the cells \(b\) and
\(c\) and leaves \(a\), \(d\) and the four marginal hypotheses unchanged.
The criterion \(a+d<b+c\) is invariant under that exchange, which proves
symmetry.

\emph{Item 2.}  Replacing \(B\) by \(\neg B\) exchanges \(a\) with \(b\)
and \(c\) with \(d\), so by item~1 \(A\) is a reason for \(\neg B\)
exactly when \(b+c<a+d\).  Replacing \(A\) by \(\neg A\) exchanges \(a\)
with \(c\) and \(b\) with \(d\), giving the same criterion \(c+b<d+a\).
By the extended-value description above,
\(\tau(B\mid A)<\tau(B\mid\neg A)\) holds in exactly the cases where
\(b+c<a+d\): the three-case analysis of item~1 applies with the inequality
reversed.  Finally, neither \(A\) nor \(\neg A\) is a reason for \(B\)
exactly when neither \(a+d<b+c\) nor \(b+c<a+d\) holds.  Since at most one
sum is infinite, this happens exactly when the two sums are equal and finite.

\emph{Item 3.}  Here \(P(A)=p_1+p_2\), \(P(B)=p_1+p_3\) and
\(p_1+p_2+p_3+p_4=1\).  Then
\[
 P(A\land B)-P(A)P(B)
 =p_1(p_1+p_2+p_3+p_4)-(p_1+p_2)(p_1+p_3)
 =p_1p_4-p_2p_3 .
\]
If \(0<P(A)<1\), the law of total probability gives
\(P(B)=P(A)P(B\mid A)+P(\neg A)P(B\mid\neg A)\), hence
\(P(B\mid A)-P(B)=P(\neg A)\bigl(P(B\mid A)-P(B\mid\neg A)\bigr)\).
The factor \(P(\neg A)\) is positive, so the two differences have the same
sign.

\emph{Item 4.}  With all cells finite, \(\tau(B\mid A)=b-a\),
\(\tau(B\mid\neg A)=d-c\) and
\(\tau(B)=\kappa(\neg B)-\kappa(B)=\min\{b,d\}-\min\{a,c\}\).
Exchanging \(A\) with \(\neg A\) exchanges the pairs \((a,b)\) and
\((c,d)\) and leaves \(\tau(B)\) unchanged, so we may assume
\(b-a\geq d-c\).  Upper bound: if \(\min\{a,c\}=a\), then
\(\tau(B)\leq b-a\) because \(\min\{b,d\}\leq b\); otherwise
\(\min\{a,c\}=c\) and \(\tau(B)\leq d-c\leq b-a\).
Lower bound: if \(\min\{b,d\}=d\), then
\(\tau(B)=d-\min\{a,c\}\geq d-c\); otherwise \(\min\{b,d\}=b\) and
\(\tau(B)=b-\min\{a,c\}\geq b-a\geq d-c\).
For \(a=0\), \(b=1\), \(c=d=5\) the marginals are \(0,5,0,1\), all
finite; \(a+d=5<6=b+c\); \(\tau(B\mid A)=1\), \(\tau(B\mid\neg A)=0\) and
\(\tau(B)=\min\{1,5\}-\min\{0,5\}=1\).  So the upper end of the interval
is attained by a genuine reason.

\emph{Item 5.}  Assigning \(0\) to every nonempty subset of \(\{1,2,3\}\)
and \(\infty\) to \(\varnothing\) satisfies normalization and the binary
minimum law, and every marginal in question is a nonempty set of rank
\(0\).  For the pair \(A=\{1\}\), \(B=\{1,2\}\) the cells are
\(\{1\},\varnothing,\{2\},\{3\}\), so \(a+d=0<\infty=b+c\).  For the pair
\(B=\{1,2\}\), \(C=\{2\}\) the cells are \(\{2\},\{1\},\varnothing,\{3\}\),
so again \(0<\infty\).  For the pair \(A=\{1\}\), \(C=\{2\}\) the cells
are \(\varnothing,\{1\},\{2\},\{3\}\), so \(a+d=\infty>0=b+c\), and by
item~2 \(A\) is a reason against \(C\).
\end{proof}
\end{annexitem}

\begin{annexitem}{annex:p15-w02b-p15-039-material}{thm:p15-w02b-p15-039-material}{Material versus conditional belief}
\begin{proof}
The negation of \(\neg A\lor B\) is \(A\land\neg B\), so the material
conditional is believed at threshold \(z\) exactly when
\(\kappa(A\land\neg B)>z\).  Since \(\kappa(A)<\infty\), the law of
conjunction of Theorem~\ref{thm:p15-rat53-conditional-laws} gives
\(\kappa(A\land\neg B)=\kappa(A)+\kappa(\neg B\mid A)\).

Item~1: \(\kappa(A)\geq0\), so
\(\kappa(A\land\neg B)\geq\kappa(\neg B\mid A)>z\).

Item~2: with \(\kappa(A)=0\) the identity reads
\(\kappa(A\land\neg B)=\kappa(\neg B\mid A)\), so the two beliefs at
threshold \(z\) are the same inequality.  For \(z=0\): by the conditional
law of negation in the same theorem, one of \(\kappa(B\mid A)\),
\(\kappa(\neg B\mid A)\) is zero.  If \(\kappa(\neg B\mid A)>0\), then
\(\kappa(B\mid A)=0\) and \(\tau(B\mid A)=\kappa(\neg B\mid A)>0\).
Conversely \(\tau(B\mid A)>0\) means
\(\kappa(\neg B\mid A)>\kappa(B\mid A)\geq0\).

Item~3: \(A\land\neg B\subseteq A\), so antitonicity
(Theorem~\ref{thm:p15-rat53-rank-laws}) gives
\(\kappa(A\land\neg B)\geq\kappa(A)>z\); the same argument with \(\neg B\)
in place of \(B\) gives belief in \(\neg A\lor\neg B\).  If both \(B\)
and \(\neg B\) were believed given \(A\) at threshold \(z\), then
\(\kappa(\neg B\mid A)>z\geq0\) and \(\kappa(B\mid A)>z\geq0\),
contradicting the conditional law of negation.

The intermediate case stated after the theorem is the displayed identity
itself.  The Tweety values used in the body are read from the prior of
Table~\ref{tab:p15-rat53-tweety}: \(\kappa(B)=0\),
\(\kappa(B\land\neg F)=\min\{2,1\}=1\), \(\kappa(P)=\min\{4,1,11,8\}=1\),
\(\kappa(P\land\neg F)=\min\{1,8\}=1\) and
\(\kappa(P\land F)=\min\{4,11\}=4\); hence \(\kappa(\neg F\mid B)=1\) and
\(\kappa(F\mid P)=3\).
\end{proof}
\end{annexitem}

\begin{annexitem}{annex:p15-w02b-p15-039-chain}{thm:p15-w02b-p15-039-chain}{The additive chain rule and minimum marginalization}
\begin{proof}
Write \(K_i=C_1\land\cdots\land C_i\).  For \(i\leq n-1\),
\(K_{n-1}\subseteq K_i\), so antitonicity gives
\(\kappa(K_i)\leq\kappa(K_{n-1})<\infty\); every conditional rank
\(\kappa(C_{i+1}\mid K_i)\) with \(i\leq n-1\) is therefore defined.
Induct on \(n\).  For \(n=2\) the claim is the law of conjunction with the
finite condition \(\kappa(C_1)\).  For \(n>2\), the law of conjunction with
the finite condition \(\kappa(K_{n-1})\) gives
\(\kappa(K_n)=\kappa(K_{n-1})+\kappa(C_n\mid K_{n-1})\), and the induction
hypothesis applies to \(C_1,\ldots,C_{n-1}\) because
\(\kappa(K_{n-2})\leq\kappa(K_{n-1})<\infty\).

For the partition, \(A=\bigcup_j(A\land D_j)\), and induction on the binary
minimum law gives \(\kappa(A)=\min_j\kappa(A\land D_j)\).  If
\(\kappa(D_j)=\infty\), then \(\kappa(A\land D_j)=\infty\) by
antitonicity, so omitting such cells does not change the minimum; at least
one cell has rank \(0\), because the minimum of the cell ranks is
\(\kappa(W)=0\).  For a finite-rank cell the law of conjunction gives
\(\kappa(A\land D_j)=\kappa(D_j)+\kappa(A\mid D_j)\).
\end{proof}
\end{annexitem}

\begin{annexitem}{annex:p15-w02b-p15-039-rational-monotony}{thm:p15-w02b-p15-039-rational-monotony}{Rational monotonicity for ranking consequence}
\begin{proof}
Not \(A\mathrel{|\!\sim}_\kappa\neg B\) means
\(\kappa(A\land B)\leq\kappa(A)\); antitonicity gives the reverse
inequality, so \(\kappa(A\land B)=\kappa(A)<\infty\), and the relation
\(A\land B\mathrel{|\!\sim}_\kappa C\) is defined.  From
\(A\mathrel{|\!\sim}_\kappa C\) and antitonicity,
\[
 \kappa(A\land B\land\neg C)\geq\kappa(A\land\neg C)>\kappa(A)
 =\kappa(A\land B),
\]
which is \(A\land B\mathrel{|\!\sim}_\kappa C\).

Tweety prior: \(\kappa(B)=0\) and \(\kappa(B\land\neg F)=1\), so
\(B\mathrel{|\!\sim}_\kappa F\).  Next \(\kappa(B\land P)=\min\{4,1\}=1\),
\(\kappa(B\land P\land F)=4>1\) and \(\kappa(B\land P\land\neg F)=1\),
so \(B\land P\mathrel{|\!\sim}_\kappa\neg F\) and not
\(B\land P\mathrel{|\!\sim}_\kappa F\): monotonicity fails.  Finally
\(\kappa(B\land P)=1>0=\kappa(B)\) is \(B\mathrel{|\!\sim}_\kappa\neg P\),
so the premise of rational monotonicity is false for this instance.
\end{proof}
\end{annexitem}
